% titel: Blatt 4, A4: Codes und Bijektivität

% Teil I: Werkzeuge, kein Lösungsweg.

% Grundidee: mod 26 rechnen; Verschiebung immer bijektiv; Multiplikation
% nur bijektiv, wenn der Faktor teilerfremd zu 26 ist (modulares
% Inverses).

Verschl\"usselung ist eine \emph{bijektive} Abbildung auf $\mathbb Z_{26}$
die 26 Buchstaben. Deshalb muss man wissen, wann eine Modulo-Abbildung
bijektiv ist. Verschiebungen ($x\mapsto x+c$) sind immer Permutationen;
Multiplikationen ($x\mapsto ax$) nur, wenn $a$ teilerfremd zu $26$ ist.

\begin{werkzeugkasten}[Rechnen mod 26]
Nach $\bmod 26$ z\"ahlt nur der Rest bei Division durch $26$; man darf
beliebig oft $26$ addieren/subtrahieren. Mit A$=0$, B$=1,\dots,$ Z$=25$:
\begin{equation}
  f(x)=(x+3)\bmod 26 .
\end{equation}
Diese Verschiebung ist bijektiv, denn auf $\mathbb Z_{26}$ ist ``$+3$``
eine feste Permutation (jeder Wert wird genau einmal getroffen); die
Umkehrung ist $f^{-1}(y)=(y-3)\bmod26$.
\end{werkzeugkasten}

\begin{werkzeugkasten}[Wann ist Multiplikation mod 26 bijektiv?]
$g(x)=ax\bmod26$ ist bijektiv genau dann, wenn $\mathrm{ggT}(a,26)=1$.
\textbf{Warum?} Bijektiv auf einer endlichen Menge $\mathbb Z_{26}$
hei\"st: injektiv (keine zwei $x$ auf denselben Wert) und surjektiv (alle
26 Reste treffen). $ax\equiv ay\ (\bmod 26)$ hei\"st $a(x-y)\equiv0$.
Ist $a$ teilerfremd zu $26$, kann man durch $a$ ``teilen`` (modulares
Inverses), also $x\equiv y$ --- injektiv, damit surjektiv. Ist
$\mathrm{ggT}(a,26)=d>1$, ist $a\cdot\frac{26}{d}\equiv0$, also z.\,B.
$x=0$ und $x=\frac{26}{d}$ kollidieren --- nicht bijektiv.
\end{werkzeugkasten}

\begin{achtung}[Typische Fehler]
\begin{itemize}
  \item Bijektiv = injektiv \emph{und} surjektiv; ein einziger nicht
        getroffener Rest (z.\,B. $1$ bei $2x$) widerlegt Surjektivität.
  \item Zwei Buchstaben mit demselben Code (Kollision): $x=0$ und
        $x=13$ bei $2x$ geben beide $0$.
  \item Die Umkehrung von $3x$ bildet man \"uber das modulare Inverse von
        $3$ (hier $9$, da $3\cdot9=27\equiv1$): multiplizieren mit $9$,
        nicht ``durch $3$`` dividieren.
\end{itemize}
\end{achtung}
